% =====================================================================
% DRAFT for the authors to reword -- main-text pieces that follow from
% moving the repeats and identifiability to the supplement (2026-09-26).
% Every statement below describes what scripts/run_noise_simulation.m
% actually does; if you change the wording, keep the facts.
% Items in [brackets] are filled from results/batches/noise_<id>/summary/.
% =====================================================================

% ---------------------------------------------------------------------
% 1. Replace the last paragraph of Section 2.4.2 ("The complete pipeline
%    was executed three times ...") with:
% ---------------------------------------------------------------------
The multi-start seed was fixed in advance at 101. Repeating the pipeline with
two further seeds (202 and 303) returned the identical accepted parameter set;
these repeats are reported in Supplementary Section~S1.

% ---------------------------------------------------------------------
% 2. Delete Section 2.4.4 (Practical identifiability). If one sentence is
%    wanted in the main text, add at the end of Section 2.4.1:
% ---------------------------------------------------------------------
The practical identifiability of the retained set at the calibrated operating
point is reported in Supplementary Section~S2.

% ---------------------------------------------------------------------
% 3. Replace Section 2.6 with:
% ---------------------------------------------------------------------
\subsection{Measurement-noise simulation}

The effect of measurement uncertainty on the calibrated parameters and on the
post-closure prediction was assessed by Monte Carlo simulation. In each of
$M = [50]$ replicates, every directly measured pre-closure target was
perturbed independently by zero-mean Gaussian noise with standard deviation
equal to its declared measurement uncertainty $\sigma_k$ (Section~2.5):
mean right atrial pressure, mean, systolic and diastolic pulmonary arterial
pressure, systolic and diastolic systemic arterial pressure, systemic flow
$Q_s$ and $Q_p/Q_s$ (Table~S[x]). Quantities that the clinical record derives
from these measurements were recomputed from the perturbed values with the same
formulas: mean arterial pressure as $(P_{\mathrm{sys}} + 2P_{\mathrm{dia}})/3$
and shunt flow as $Q_s(Q_p/Q_s - 1)$. Draws that were physically impossible
(a non-positive value, diastolic above systolic pressure, mean pulmonary
pressure outside its systolic--diastolic range, or $Q_p/Q_s \le 1$) were
redrawn; [n] draws were redrawn in total. Model inputs (age, weight, height,
heart rate and defect diameter) and all post-closure measurements were not
perturbed. The noise draws used a fixed seed (20260926).

Each perturbed record was recalibrated with the procedure of Section~2.4.2
unchanged: the same clinical seeding, staged objective, six starting points
and multi-start seed (101). The fitted parameter set was held at the seven
parameters selected by the sensitivity screen of Section~2.4.1 and was not
re-screened, so that only the measurements differ between replicates. The
recalibrated parameters were then applied with the defect closed, as in
Section~2.4.3, and the prediction was compared with the unperturbed
post-closure measurements, which remained withheld from calibration.

For each calibrated parameter and each predicted quantity we report the mean,
standard deviation, coefficient of variation and the 2.5th--97.5th percentile
interval across replicates, and for each post-closure target the fraction of
replicates within the $\pm10\%$ gate and whether the measured value lies
inside the 95\% interval of the prediction. All replicates are reported,
including any that failed or were not accepted.

% ---------------------------------------------------------------------
% 4. Section 3.7 -- replace the current table ENTIRELY. The current table
%    (Rep 1-3, RMSE 0.0627 / 0.0546 / 0.0822, measured MAP 77) is not a
%    noise simulation: it is the superseded 12-parameter three-seed batch.
%    Suggested structure once the batch has run:
% ---------------------------------------------------------------------
\subsection{Measurement-noise simulation}

[M] of [M] replicates completed; [k] were classified accepted, [k]
physiological but poor fit and [k] not accepted. Across replicates the
governed RMSE against the perturbed record was [median] (range [a]--[b]).

% Table A -- calibrated parameters: noise_summary_parameters.csv
%   Parameter | Reference (Table 3) | Mean +- SD | CV (%) | 95% interval
% Table B -- post-closure prediction: noise_summary_post.csv
%   Target | Measured | Predicted (Table 7) | Mean [95% interval] |
%   within 10% (% of replicates) | measured inside interval
